% Worked Examples: Topic 3.7 - PID Control
\documentclass[11pt,a4paper]{article}
\usepackage{../worked-examples-template}

\title{\textbf{Worked Examples}\\[0.5em]
\large Topic 3.7: PID Control\\[0.3em]
\normalsize \coursesubtitle{} -- \coursetitle}
\author{Assoc.\ Prof.\ William Robertson \and Dr Sean McGowan}
\date{\today}

\begin{document}

\maketitle
\tableofcontents
\newpage

\section{Concept 3.7.1: Feedback Control}

\subsection{Example 1: Basic Feedback Control Design}

\paragraph{Problem:} A temperature control system has the plant transfer function:
\begin{align}
P(s) = \frac{2}{s+1}
\end{align}

Design a proportional controller $C(s) = K_p$ for a unity feedback configuration such that the closed-loop system has:
\begin{itemize}
\item Steady-state error less than 10\% for a unit step input
\item Settling time (2\% criterion) less than 4 seconds
\end{itemize}

\paragraph{Solution:}

For unity feedback with controller $C(s) = K_p$ and plant $P(s)$:

The closed-loop transfer function is:
\begin{align}
T(s) = \frac{K_pP(s)}{1 + K_pP(s)} = \frac{2K_p}{s + 1 + 2K_p}
\end{align}

\textbf{Requirement 1: Steady-State Error}

For a unit step input, the steady-state error is:
\begin{align}
e_{ss} = \frac{1}{1 + K_pP(0)} = \frac{1}{1 + 2K_p}
\end{align}

For $e_{ss} < 0.1$ (10\%):
\begin{align}
\frac{1}{1 + 2K_p} < 0.1 \quad \Rightarrow \quad 1 + 2K_p > 10 \quad \Rightarrow \quad K_p > 4.5
\end{align}

\textbf{Requirement 2: Settling Time}

The closed-loop pole is at:
\begin{align}
s = -(1 + 2K_p)
\end{align}

For a first-order system, the settling time (2\% criterion) is:
\begin{align}
T_s = \frac{4}{|s|} = \frac{4}{1 + 2K_p}
\end{align}

For $T_s < 4$ seconds:
\begin{align}
\frac{4}{1 + 2K_p} < 4 \quad \Rightarrow \quad 1 + 2K_p > 1 \quad \Rightarrow \quad K_p > 0
\end{align}

This is always satisfied for positive $K_p$.

\textbf{Design Choice:}

The binding constraint is $K_p > 4.5$. Choose $K_p = 5$ for some margin.

With $K_p = 5$:
\begin{itemize}
\item Closed-loop transfer function: $T(s) = \frac{10}{s+11}$
\item Steady-state error: $e_{ss} = \frac{1}{1+10} = 0.091 = 9.1\%$ ✓
\item Settling time: $T_s = \frac{4}{11} \approx 0.36$ seconds ✓
\end{itemize}

\textbf{Result:} The proportional controller $C(s) = 5$ meets both design requirements.

\section{Concept 3.7.2: Simple Controllers}

\subsection{Example 1: Comparing P, PI, and PD Controllers}

\paragraph{Problem:} For the plant:
\begin{align}
P(s) = \frac{1}{s(s+2)}
\end{align}

Compare the closed-loop performance with three different controllers:
\begin{itemize}
\item[(a)] Proportional (P): $C(s) = K_p = 4$
\item[(b)] Proportional-Integral (PI): $C(s) = K_p + \frac{K_i}{s} = 4 + \frac{2}{s}$
\item[(c)] Proportional-Derivative (PD): $C(s) = K_p + K_d s = 4 + 2s$
\end{itemize}

For each controller, find the closed-loop transfer function and steady-state error to a unit step input.

\paragraph{Solution:}

\textbf{Part (a): Proportional Controller}

\begin{align}
T_P(s) = \frac{C(s)P(s)}{1 + C(s)P(s)} = \frac{4/[s(s+2)]}{1 + 4/[s(s+2)]} = \frac{4}{s^2 + 2s + 4}
\end{align}

For a Type-1 system (one integrator in the loop), the steady-state error to a step is:
\begin{align}
e_{ss} = \lim_{s \to 0} \frac{sR(s)}{1 + C(s)P(s)} = \lim_{s \to 0} \frac{1}{1 + 4/[s(s+2)]}
\end{align}

As $s \to 0$, the denominator goes to infinity, so $e_{ss} = 0$.

\textbf{Part (b): PI Controller}

\begin{align}
C(s) = 4 + \frac{2}{s} = \frac{4s + 2}{s}
\end{align}

The loop transfer function:
\begin{align}
C(s)P(s) = \frac{4s+2}{s} \cdot \frac{1}{s(s+2)} = \frac{4s+2}{s^2(s+2)}
\end{align}

Closed-loop transfer function:
\begin{align}
T_{PI}(s) = \frac{4s+2}{s^2(s+2) + 4s + 2} = \frac{4s+2}{s^3 + 2s^2 + 4s + 2}
\end{align}

For steady-state error to step:
\begin{align}
e_{ss} = \lim_{s \to 0} \frac{s}{1 + C(s)P(s)} = \lim_{s \to 0} \frac{s^2(s+2)}{s^2(s+2) + 4s + 2} = \frac{0}{2} = 0
\end{align}

\textbf{Part (c): PD Controller}

\begin{align}
C(s) = 4 + 2s
\end{align}

The loop transfer function:
\begin{align}
C(s)P(s) = (4 + 2s) \cdot \frac{1}{s(s+2)} = \frac{2s+4}{s(s+2)}
\end{align}

Closed-loop transfer function:
\begin{align}
T_{PD}(s) = \frac{2s+4}{s(s+2) + 2s + 4} = \frac{2s+4}{s^2 + 4s + 4} = \frac{2s+4}{(s+2)^2}
\end{align}

For steady-state error:
\begin{align}
e_{ss} = \lim_{s \to 0} \frac{s}{1 + C(s)P(s)} = \lim_{s \to 0} \frac{s(s+2)}{s(s+2) + 2s + 4} = \frac{0}{4} = 0
\end{align}

\textbf{Summary Comparison:}

\begin{center}
\begin{tabular}{|c|c|c|}
\hline
Controller & Closed-Loop Poles & $e_{ss}$ (step) \\
\hline
P: $K_p = 4$ & $s = -1 \pm j\sqrt{3}$ & 0 \\
PI: $K_p = 4, K_i = 2$ & $s^3 + 2s^2 + 4s + 2 = 0$ & 0 \\
PD: $K_p = 4, K_d = 2$ & $s = -2$ (double) & 0 \\
\hline
\end{tabular}
\end{center}

\textbf{Observations:}
\begin{itemize}
\item All achieve zero steady-state error (plant has integrator)
\item PD gives critically damped response (fastest settling, no overshoot)
\item PI adds extra pole, making system third-order (may be slower)
\item P gives oscillatory response with some overshoot
\end{itemize}

\section{Concept 3.7.3: PID Control}

\subsection{Example 1: PID Controller Tuning using Ziegler-Nichols}

\paragraph{Problem:} A process has the First-Order Time-Delay (FOTD) model:
\begin{align}
P(s) = \frac{2}{1+5s}e^{-s}
\end{align}

Use the Ziegler-Nichols time-domain method to design a PID controller for this system. The parameters are $K = 2$, $T = 5$, and $\tau = 1$.

\paragraph{Solution:}

For the FOTD model $P(s) = \frac{K}{1+Ts}e^{-\tau s}$, the Ziegler-Nichols method uses the ITD approximation with:
\begin{itemize}
\item Slope parameter: $a = \frac{K}{T} = \frac{2}{5} = 0.4$
\item Time delay: $\tau = 1$
\end{itemize}

The standard PID controller form is:
\begin{align}
C(s) = K_p\left(1 + \frac{1}{T_i s} + T_d s\right)
\end{align}

According to Ziegler-Nichols tuning rules for PID control (time-domain method):
\begin{align}
K_p &= \frac{1.2}{a\tau} = \frac{1.2}{0.4 \times 1} = 3 \\
T_i &= 2\tau = 2 \times 1 = 2 \\
T_d &= 0.5\tau = 0.5 \times 1 = 0.5
\end{align}

Therefore, the PID controller is:
\begin{align}
C(s) = 3\left(1 + \frac{1}{2s} + 0.5s\right) = 3 + \frac{3}{2s} + 1.5s
\end{align}

In standard form:
\begin{align}
C(s) = K_p + \frac{K_i}{s} + K_d s
\end{align}

where:
\begin{itemize}
\item Proportional gain: $K_p = 3$
\item Integral gain: $K_i = \frac{K_p}{T_i} = \frac{3}{2} = 1.5$
\item Derivative gain: $K_d = K_p T_d = 3 \times 0.5 = 1.5$
\end{itemize}

\textbf{Controller Transfer Function:}
\begin{align}
C(s) = 3 + \frac{1.5}{s} + 1.5s = \frac{1.5s^2 + 3s + 1.5}{s}
\end{align}

\textbf{Expected Performance:}
\begin{itemize}
\item The Ziegler-Nichols method typically gives $\approx$ 25\% overshoot
\item Provides good disturbance rejection
\item May need fine-tuning for specific performance requirements
\end{itemize}

\textbf{Practical Implementation Note:}

In practice, the derivative term is often implemented with a low-pass filter to avoid noise amplification:
\begin{align}
C(s) = K_p + \frac{K_i}{s} + \frac{K_d s}{1 + \frac{T_d s}{N}}
\end{align}

where $N = 10$ is typical (giving filter pole at $s = -10/T_d = -20$).

\section{Concept 3.7.4: Anti-Windup}

\subsection{Example 1: Integrator Windup Problem and Solution}

\paragraph{Problem:} A PI controller is used to control a motor with position output. The controller is:
\begin{align}
u(t) = K_p e(t) + K_i \int_0^t e(\tau) d\tau
\end{align}

with $K_p = 5$ and $K_i = 2$. The control signal is limited to $|u| \leq 10$ volts.

(a) Explain what happens when a large step reference is applied that causes $u(t)$ to saturate.

(b) Design an anti-windup scheme using back-calculation with time constant $T_t = 1$ second.

\paragraph{Solution:}

\textbf{Part (a): Windup Problem}

When a large step reference is applied:

\begin{enumerate}
\item The error $e(t) = r(t) - y(t)$ is initially large
\item The controller output $u(t)$ quickly exceeds the limit and saturates at $u = 10$ V
\item While $u$ is saturated, the error remains large (output hasn't caught up to reference)
\item The integral term continues to accumulate: $\int e(\tau) d\tau$ keeps growing
\item This is called \textbf{integrator windup}
\item When the output finally approaches the reference, the integral has become very large
\item The controller output stays saturated even after the error becomes negative
\item This causes large overshoot and poor transient response
\end{enumerate}

\textbf{Part (b): Anti-Windup with Back-Calculation}

The anti-windup scheme adds a feedback path that prevents the integrator from winding up:

\begin{align}
\frac{dx}{dt} = K_i e - \frac{1}{T_t}(u_{sat} - u)
\end{align}

where:
\begin{itemize}
\item $x$ is the integral state
\item $u = K_p e + x$ is the controller output before saturation
\item $u_{sat} = \text{sat}(u)$ is the saturated control signal
\item $T_t = 1$ s is the tracking time constant
\end{itemize}

The saturated control is:
\begin{align}
u_{sat} = \begin{cases}
10 & \text{if } u > 10 \\
u & \text{if } -10 \leq u \leq 10 \\
-10 & \text{if } u < -10
\end{cases}
\end{align}

\textbf{How it works:}

\begin{itemize}
\item When $u$ is within limits: $u_{sat} = u$, so $\frac{dx}{dt} = K_i e$ (normal operation)
\item When $u$ saturates: The difference $(u_{sat} - u)$ is non-zero
\item This difference reduces the integration rate: $\frac{dx}{dt} = K_i e - \frac{1}{T_t}(u_{sat} - u)$
\item The integrator state is "tracked back" to a value consistent with the saturation limit
\item The time constant $T_t$ determines how quickly the tracking occurs
\item With $T_t = 1$ s, the integrator adjusts relatively quickly
\end{itemize}

\textbf{Implementation:}

The complete controller with anti-windup is:
\begin{align}
u &= K_p e + x \\
u_{sat} &= \text{sat}(u, -10, 10) \\
\frac{dx}{dt} &= K_i e - \frac{1}{1}(u_{sat} - u) = 2e - (u_{sat} - u)
\end{align}

\textbf{Benefits:}
\begin{itemize}
\item Prevents large overshoot caused by integrator windup
\item Maintains good tracking during saturation
\item Minimal impact on performance when not saturating
\item Simple to implement in both analog and digital systems
\end{itemize}

\section{Concept 3.7.5: Implementation}

\subsection{Example 1: Discrete-Time PID Implementation}

\paragraph{Problem:} A continuous PID controller:
\begin{align}
u(t) = K_p e(t) + K_i \int_0^t e(\tau)d\tau + K_d \frac{de(t)}{dt}
\end{align}

with $K_p = 2$, $K_i = 0.5$, $K_d = 1$ needs to be implemented digitally with sampling period $T_s = 0.1$ seconds.

(a) Derive the discrete-time implementation using:
\begin{itemize}
\item Forward Euler for the integral
\item Backward difference for the derivative
\end{itemize}

(b) Write the resulting difference equation.

\paragraph{Solution:}

\textbf{Part (a): Discretization}

Let $e_k = e(kT_s)$ denote the error at time step $k$.

\textbf{Proportional term:} No approximation needed
\begin{align}
u_P[k] = K_p e_k = 2e_k
\end{align}

\textbf{Integral term:} Using forward Euler (rectangular rule)
\begin{align}
\int_0^{kT_s} e(\tau)d\tau \approx T_s \sum_{j=0}^{k} e_j
\end{align}

Define the accumulated sum:
\begin{align}
I_k = I_{k-1} + T_s e_k
\end{align}

Then:
\begin{align}
u_I[k] = K_i I_k = 0.5 I_k
\end{align}

\textbf{Derivative term:} Using backward difference
\begin{align}
\frac{de}{dt}\bigg|_{t=kT_s} \approx \frac{e_k - e_{k-1}}{T_s}
\end{align}

Therefore:
\begin{align}
u_D[k] = K_d \frac{e_k - e_{k-1}}{T_s} = 1 \cdot \frac{e_k - e_{k-1}}{0.1} = 10(e_k - e_{k-1})
\end{align}

\textbf{Part (b): Complete Difference Equation}

The total control signal is:
\begin{align}
u_k = u_P[k] + u_I[k] + u_D[k]
\end{align}

Substituting:
\begin{align}
u_k = 2e_k + 0.5I_k + 10(e_k - e_{k-1})
\end{align}

With the integral update:
\begin{align}
I_k = I_{k-1} + 0.1e_k
\end{align}

Combining:
\begin{align}
u_k = 2e_k + 0.5(I_{k-1} + 0.1e_k) + 10e_k - 10e_{k-1}
\end{align}

Simplifying:
\begin{align}
u_k = 0.5I_{k-1} + (2 + 0.05 + 10)e_k - 10e_{k-1}
\end{align}

\begin{align}
u_k = 0.5I_{k-1} + 12.05e_k - 10e_{k-1}
\end{align}

\textbf{Algorithm:}

\begin{verbatim}
Initialize: I = 0, e_prev = 0

At each time step k:
1. Read measurement y_k
2. Compute error: e_k = r_k - y_k
3. Compute control:
   u_k = 0.5*I + 12.05*e_k - 10*e_prev
4. Update integral: I = I + 0.1*e_k
5. Update previous error: e_prev = e_k
6. Apply control signal: u_k (with saturation if needed)
\end{verbatim}

\textbf{Alternative Velocity Form:}

For better numerical properties and easier anti-windup:
\begin{align}
\Delta u_k = u_k - u_{k-1} = q_0 e_k + q_1 e_{k-1} + q_2 e_{k-2}
\end{align}

where:
\begin{align}
q_0 &= K_p + K_i T_s + \frac{K_d}{T_s} = 2 + 0.05 + 10 = 12.05 \\
q_1 &= -K_p - 2\frac{K_d}{T_s} = -2 - 20 = -22 \\
q_2 &= \frac{K_d}{T_s} = 10
\end{align}

This form computes the change in control signal:
\begin{align}
u_k = u_{k-1} + 12.05e_k - 22e_{k-1} + 10e_{k-2}
\end{align}

\textbf{Practical Considerations:}
\begin{itemize}
\item Add low-pass filter on derivative to reduce noise sensitivity
\item Implement anti-windup for integral term
\item Use appropriate data types (floating-point for better precision)
\item Consider derivative on measurement rather than error to avoid derivative kick
\end{itemize}

\end{document}
